\documentclass[11pt,a4paper]{article}

\usepackage[margin=25mm]{geometry}
\usepackage[T1]{fontenc}
\usepackage[utf8]{inputenc}
\usepackage{lmodern}
\usepackage{microtype}
\usepackage{amsmath,amssymb}
\usepackage{booktabs}
\usepackage{tabularx}
\usepackage{array}
\usepackage{graphicx}
\usepackage{xcolor}
\usepackage{enumitem}
\usepackage{hyperref}
\usepackage[nameinlink,capitalise]{cleveref}
\usepackage{csquotes}

\hypersetup{colorlinks=true, linkcolor=blue!55!black, citecolor=blue!55!black, urlcolor=blue!55!black}
\setlength{\parindent}{0pt}
\setlength{\parskip}{0.55em}
\renewcommand{\arraystretch}{1.18}
\newcolumntype{Y}{>{\raggedright\arraybackslash}X}
\newcommand{\always}{\mathop{\square}}
\newcommand{\eventually}{\mathop{\Diamond}}
\newcommand{\ruleid}[1]{\texttt{#1}}

\title{\textbf{ATIC: A Formal Rulebook for Autonomous Drone Operations}\\
\large From Regulation to Specification}
\author{Mamoun Rhmari Tlemcani \and Hadrien Legros \and Ben Youssef Belkis \and Saad Lahlou\\[0.4em]
\normalsize ETH Z\"urich --- Advanced Topics in Control (ATIC)}
\date{July 2026}

\begin{document}
\maketitle

\begin{abstract}
Autonomous drones must satisfy rules that are written for people, operators, manufacturers, and public authorities rather than for a flight controller. We construct a machine-readable rulebook that translates the runtime-relevant parts of the European UAS and U-space regulations into Signal Temporal Logic (STL)-style requirements and orders them by priority. The resulting rulebook contains applicability gates and eight operational priority tiers, from human safety to mission progress. We deliberately separate these behavioural rules from registration, certification, and other administrative obligations that cannot be decided from a flight trajectory alone. A lightweight Python monitor evaluates simulated trajectories using signed robustness margins. Three synthetic case studies demonstrate (i) route re-planning around people and a quiet zone, (ii) mission abortion and safe landing under a low-battery constraint, and (iii) reconfiguration after an emergency area becomes active during the flight. In each case, the compliant alternative satisfies all monitored rules at all 360 samples, whereas the naive trajectory achieves only 26.4\%, 40.6\%, and 60.3\% joint compliance, respectively. The work is a formalisation and monitoring prototype, not a legal compliance certificate or a map of local Z\"urich restrictions.
\end{abstract}

\section{Problem formulation and scope}

The project studies a single autonomous unmanned aircraft performing an urban or semi-urban delivery or inspection mission. Its objective is not merely to find a short route from a start point to a destination. The aircraft must also respect legal restrictions such as altitude ceilings, geographical zones, proximity to people, vehicle-operation limits, and short-notice airspace changes.

The difficulty is twofold. First, regulations are expressed in natural language, with exceptions, conditions, and terms that are not directly executable. Second, the rules do not have equal standing: avoiding an assembly of people must not be traded for a shorter delivery route. Our goal is therefore to answer the following questions:

\begin{enumerate}[leftmargin=1.4em]
    \item How can relevant drone regulations be represented as unambiguous, monitorable temporal specifications over a flight trace?
    \item How can the rules be organised so that safety and legal constraints take precedence over operational convenience and mission completion?
\end{enumerate}

The output is a formal rulebook and a monitoring demonstrator. Given a trajectory $x(t)$, the monitor evaluates each selected rule, reports a Boolean verdict and a signed robustness margin, and visualises the first violated rule in priority order. A trajectory optimiser then produces an alternative path with lower aggregate violation depth and a shorter path where this is compatible with the monitored constraints.

\subsection{Regulatory scope}

Our source review is based on the supplied EASA \emph{Easy Access Rules for Unmanned Aircraft Systems} (Regulations (EU) 2019/947 and 2019/945) and \emph{Easy Access Rules for U-space} (Regulation (EU) 2021/664). Switzerland applies these EU UAS regulations through the EU--Switzerland aviation agreement; the Swiss FOCA acts as the competent authority. The supplied material does \emph{not} contain a separate FOCA geographical-zone data set or a Z\"urich-specific ordinance. Accordingly, this report formalises European and Swiss-applicable rules in general, not special rules for Z\"urich.

The maps used in the experiments are intentionally synthetic. Labels such as ``hospital helipad'' and ``emergency response area'' illustrate the geometry and temporal activation of a rule, but do not assert the existence, extent, or legal status of a real restricted area. The separate city visualisation in the repository uses OpenStreetMap building geometry only as a presentation aid.

An autonomous operation also cannot simply be treated as an \emph{open}-category flight. The rulebook captures the gate $\always(\texttt{autonomous}\Rightarrow \neg\texttt{open})$: a genuinely autonomous mission must be assessed in the specific or certified branch and requires an appropriate legal basis. This is a classification constraint before take-off, rather than a trade-off that a trajectory optimiser may resolve during flight.

\subsection{What is intentionally out of scope}

The regulations also specify operator registration, physical labels, manufacturer conformity procedures, training records, competent-authority duties, certification, and market surveillance. These are important legal obligations, but most are not runtime properties of a trajectory. We retain them in the natural-language inventory for traceability, while excluding purely administrative, organisational, manufacturing, and labelling obligations from the trajectory-level STL monitor. A future operational system would need a separate pre-flight compliance module for these gates.

\section{From regulations to a prioritised rulebook}

\subsection{Formalisation approach}

We model a mission as a time-indexed signal containing position $p(t)=(x(t),y(t),z(t))$, speed $v(t)$, height above ground level $h_{\mathrm{AGL}}(t)$, battery level $b(t)$, and discrete context signals such as \texttt{night}, \texttt{emerg\_area}, and \texttt{light\_on}. The operators $\always$ and $\eventually$ denote ``always'' and ``eventually'', respectively. For example, the 120~m ceiling is expressed as
\begin{equation}
    \varphi_{\mathrm{alt}} = \always_{[0,T]}\bigl(h_{\mathrm{AGL}}(t) \leq c(t)\bigr),
    \label{eq:altitude}
\end{equation}
where $c(t)$ is normally 120~m but may be changed by applicable zone data or an explicit legal exception. A simple spatial exclusion rule for a circular zone with centre $q$ and radius $r$ is
\begin{equation}
    \varphi_{\mathrm{zone}} = \always_{[0,T]}\bigl(\lVert p_{xy}(t)-q\rVert > r\bigr).
    \label{eq:zone}
\end{equation}

For monitoring and optimisation, each predicate also receives a signed robustness. For example, $\rho_{\mathrm{zone}}(t)=\lVert p_{xy}(t)-q\rVert-r$; positive values are a spatial safety margin, zero is the boundary, and negative values quantify penetration. The monitor evaluates the relevant predicates at 360 equally spaced samples over a mission. It reports a rule as passing only when its robustness remains non-negative throughout the trace.

The initial natural-language inventory contains 37+ requirements with exact legal sources, triggers, exceptions, and notes. We organise the executable requirements into the rulebook in \cref{tab:tiers}. Tier~0 selects the applicable legal branch. Tiers~1--7 are ordered constraints; Tier~8 is a progress objective considered only after the constraints. Rules within one tier are intentionally incomparable unless a legal source gives a more precise order.

\begin{table}[htbp]
\centering
\caption{Priority structure of the rulebook. Lower tier number means higher priority.}
\label{tab:tiers}
\small
\begin{tabularx}{\linewidth}{@{}c l Y@{}}
\toprule
Tier & Theme & Examples \\
\midrule
0 & Applicability gates & Category determination; autonomy is not permitted in the open category. \\
1 & Safety of humans & No overflight of assemblies; hazard avoidance; emergency-area exclusion; flight termination. \\
2 & Containment and property safety & Operational-volume containment; geofence behaviour; dangerous-goods restrictions. \\
3 & Airspace and geographical zones & 120~m AGL ceiling; no-fly zones; U-space authorisation and dynamic reconfiguration. \\
4 & Identification and registration & Registration and remote-identification checks. \\
5 & Vehicle control and operating limits & C2-link handling; low-battery landing; speed and payload limits. \\
6 & Conspicuity & Night lighting. \\
7 & Comfort & Noise limit. \\
8 & Mission progress & Reach the goal while minimising mission time/path length. \\
\bottomrule
\end{tabularx}
\end{table}

\subsection{Rules selected for the demonstrations}

The three case studies use the subset in \cref{tab:selected-rules}. This is a deliberate selection, not a claim that these are the only applicable rules. In particular, registration and authorisation should be checked before real deployment even when they are not visualised on a map.

\begin{table}[htbp]
\centering
\caption{Selected runtime rules and their monitorable abstraction.}
\label{tab:selected-rules}
\small
\begin{tabularx}{\linewidth}{@{}l c Y Y@{}}
\toprule
Rule & Tier & Runtime interpretation & STL-style form \\
\midrule
\ruleid{REQ-PROX-01} & 1 & Do not overfly an assembly of people. & $\always\neg\texttt{assembly\_below}$ \\
\ruleid{REQ-PLT-03} & 1 & Divert if an emergency area is active and no permission is present. & $\always(\texttt{emerg\_area}\wedge\neg\texttt{permission}\Rightarrow\eventually_{[0,\tau]}\texttt{diverted})$ \\
\ruleid{REQ-PROX-06} & 2 & Remain within the declared operational volume. & $\always\texttt{in\_op\_vol}$ \\
\ruleid{REQ-ZONE-03} & 3 & Remain at or below the applicable AGL ceiling. & $\always(h_{\mathrm{AGL}}\leq c)$ \\
\ruleid{REQ-ZONE-04} & 3 & Satisfy each applicable geographical-zone constraint. & $\always(\texttt{in\_zone}_k\Rightarrow\texttt{satisfies}(k))$ \\
\ruleid{REQ-USP-08} & 3 & Become compliant after a short-notice airspace reconfiguration. & $\always(\texttt{reconfig\_notice}\Rightarrow\eventually_{[0,5s]}\texttt{reconfig\_compliant})$ \\
\ruleid{REQ-EQP-03} & 5 & On low battery, warn and land with adequate lead time. & $\always(b\leq b_\mathrm{warn}\Rightarrow\eventually_{[0,t_\mathrm{land}]}\texttt{landed})$ \\
\ruleid{REQ-EQP-07} & 5 & C1 maximum speed is 19~m/s. & $\always(v\leq19)$ \\
\ruleid{REQ-EQP-04} & 6 & Operate with the required light at night. & $\always(\texttt{night}\Rightarrow\texttt{light\_on})$ \\
\ruleid{REQ-EQP-12} & 7 & Respect the applicable noise limit. & $\always(\texttt{noise}\leq\texttt{limit})$ \\
\bottomrule
\end{tabularx}
\end{table}

\subsection{Monitoring and trajectory generation}

The code in \texttt{scenario/core.py} implements a common scenario interface, geometric signed-distance monitors, trajectory sampling, a compliance report, and an optimiser. Each case specifies a start point, a delivery point (or a safe landing point), zones, a naive route, and a fallback compliant route. A Catmull--Rom interpolation creates a smooth trace through waypoints. The alternative trajectory is obtained either by L-BFGS-B optimisation of path length plus aggregate violation depth or by using the handcrafted fallback path when it has lower violation depth.

The monitor uses the priority order to display the highest-priority currently violated rule. It does \emph{not yet} implement a fully lexicographic, tier-by-tier optimisation objective: the numerical penalty sums violation depths across selected rules. The hierarchy is therefore explicit in the formal rulebook and in verdict reporting, but strict priority preservation in synthesis remains future work. This distinction is important: an implementation should not claim that a weighted aggregate objective is equivalent to the rulebook's lexicographic order.

\section{Scenario experiments and results}

Each experiment compares a direct, rule-unaware trajectory with an alternative trajectory under the same synthetic environment. The green path in the figures is the monitored compliant path; the red dashed path is the naive reference. All results below were obtained with 360 samples per mission.

\subsection{Case 1: standard route re-planning}

The first mission crosses two crowd zones and a residential quiet zone, and continues from daytime into nighttime. The direct route is short but cuts through an assembly of people and the quiet zone. The compliant route goes around these areas while respecting the altitude and speed limits. The light is enabled when the night signal becomes active; this is a configuration change in addition to the primary spatial route change.

\begin{figure}[htbp]
\centering
\includegraphics[width=0.93\linewidth]{figures/crowd_noise_night.png}
\caption{Case 1: a spatial reroute avoids crowd and residential quiet zones.}
\label{fig:crowd}
\end{figure}

The naive trace violates crowd standoff for 73 samples, the night-light requirement for 170 samples, and the quiet-zone noise requirement for 83 samples. Its joint compliance is 26.4\%. The compliant route passes all five selected rules over the full mission, including a minimum 80.4~m crowd-standoff margin.

\subsection{Case 2: low battery changes the mission objective}

The second case demonstrates that delivery completion is the lowest-priority objective. A restricted zone lies between the start and the original delivery point. Battery decreases from 45\% to 8\%, with a warning threshold of 20\%. Once the low-battery condition is relevant, the correct behaviour is not to force completion of the delivery. The compliant trajectory avoids the restricted zone, heads to the designated safe landing site, descends to ground level, and holds there for the remaining simulated horizon.

\begin{figure}[htbp]
\centering
\includegraphics[width=0.93\linewidth]{figures/battery_abort.png}
\caption{Case 2: safe landing takes precedence over reaching the original delivery point.}
\label{fig:battery}
\end{figure}

The direct trajectory enters the restricted zone for 97 samples and has no valid low-battery landing response for 117 samples, resulting in 40.6\% joint compliance. The safe-landing route is 100\% compliant and retains a minimum low-battery landing robustness of $+4.8$ in the simplified landing model. This is the key qualitative result: completing the package-delivery objective is intentionally abandoned to satisfy a higher-priority operational safety rule.

\subsection{Case 3: dynamic emergency-area activation}

The third case separates static and dynamic constraints. It contains a permanent hospital-helipad no-fly zone and an emergency response area that becomes active 65~s after take-off. The naive trajectory crosses both areas and climbs above the 120~m limit. The alternative reroutes along the boundary of the operational volume after activation, thereby maintaining clearance from both the static and dynamically activated zones.

\begin{figure}[htbp]
\centering
\includegraphics[width=0.93\linewidth]{figures/emergency_reconfiguration.png}
\caption{Case 3: a temporary emergency area activates during flight; the alternative maintains the required exclusion.}
\label{fig:emergency}
\end{figure}

The naive trajectory violates emergency-area exclusion for 46 samples, the static no-fly zone for 63 samples, and the altitude ceiling for 97 samples. Despite staying inside the broad operational volume and below the speed cap, its joint compliance is only 60.3\%. The monitored alternative passes every selected rule at all 360 samples. It maintains a $+122.1$~m emergency-zone margin and a $+168.0$~m static no-fly-zone margin in the synthetic geometry.

\begin{table}[htbp]
\centering
\caption{Joint monitoring results. A pass requires every selected rule to be satisfied at every sample.}
\label{tab:results}
\small
\begin{tabularx}{\linewidth}{@{}Y c c Y@{}}
\toprule
Case & Naive joint compliance & Alternative joint compliance & Main outcome \\
\midrule
Route re-planning & 95/360 (26.4\%) & 360/360 (100\%) & Reroute avoids crowds and quiet zone; night light is enabled. \\
Low-battery abort & 146/360 (40.6\%) & 360/360 (100\%) & Mission ends at the safe landing site rather than the original delivery point. \\
Dynamic emergency area & 217/360 (60.3\%) & 360/360 (100\%) & Route remains clear after the emergency area activates at 65~s. \\
\bottomrule
\end{tabularx}
\end{table}

\section{Discussion}

The experiments demonstrate that a formal rulebook is useful even with simple geometry. It makes a rule violation explicit, quantifies its magnitude, and exposes when goal completion should be sacrificed. The battery case is particularly instructive: a trajectory that reaches the delivery location is not necessarily the correct trajectory once a higher-priority landing obligation becomes active. Likewise, the emergency scenario shows why the controller needs access to dynamic airspace information rather than only a fixed map loaded before take-off.

Several limitations delimit the result. First, the experiments abstract people, airspace, noise, and landing sites as circles and simple signals. They do not model wind, perception error, communication delay, obstacle avoidance, detailed 3D terrain, or aircraft dynamics. In particular, the 120~m rule is represented as a fixed ceiling in the scenarios, whereas the regulation refers to height above the closest ground surface and permits narrowly defined exceptions. Second, some formal rules need inputs that are external to the controller, such as the validity of an authorisation, live U-space information, or emergency-service permission. The report formalises the response once those signals are supplied; it does not create those services.

Third, the rule inventory is broader than the visual monitor. Administrative obligations are not less important; they require a different kind of verification. Finally, although the rulebook provides a priority ordering, the current optimisation implementation uses a summed penalty rather than a strict lexicographic solver. A natural next step is hierarchical optimisation: minimise Tier~1 violation first, then Tier~2 without worsening Tier~1, and so on, before optimising path length. This would align synthesis exactly with the rulebook semantics described by Censi et al.

\section{Conclusion}

We translated a substantial set of EASA UAS and U-space provisions into a prioritised, monitorable rulebook for autonomous-drone operation. The work keeps the distinction between legal applicability gates, runtime flight constraints, and non-trajectory administrative obligations. The three scenario studies show route re-planning, mission abortion for safe landing, and response to a dynamic exclusion area. All alternative trajectories satisfy the selected monitored rules, while their naive counterparts do not.

The report should be read as a formal-methods prototype: it demonstrates how natural-language regulation can drive monitorable controller requirements and expose priority-sensitive mission decisions. It is not a substitute for obtaining current FOCA zone data, a flight authorisation, or a full safety assessment before a real operation.

\section*{Reproducibility}

The implementation is included with this report. Run the following command from the repository root to generate the three interactive scenario visualisations and SVG maps:
\begin{center}
\texttt{python scenario/scenario\_simulation.py}
\end{center}
The relevant modules are \texttt{scenario/core.py}, \texttt{scenario/scene.py}, and \texttt{scenario/scenarios/}. The natural-language inventory and formal rulebook are provided as \texttt{Rules\_Plain\_Language.pdf} and \texttt{STL\_Rules.pdf}.

\begin{thebibliography}{9}

\bibitem{censi2019}
A. Censi et al.
\newblock \emph{Liability, Ethics, and Culture-Aware Behavior Specification Using Rulebooks}.
\newblock In \emph{Proceedings of the IEEE International Conference on Robotics and Automation (ICRA)}, 2019.

\bibitem{rizaldi2017}
A. Rizaldi et al.
\newblock \emph{Formalising and Monitoring Traffic Rules for Autonomous Vehicles in Isabelle/HOL}.
\newblock In \emph{Integrated Formal Methods (iFM)}, 2017.

\bibitem{easaUAS}
European Union Aviation Safety Agency.
\newblock \emph{Easy Access Rules for Unmanned Aircraft Systems: Regulations (EU) 2019/947 and 2019/945}.
\newblock Supplied project source, June 2026 edition.

\bibitem{easaUspace}
European Union Aviation Safety Agency.
\newblock \emph{Easy Access Rules for U-space: Regulation (EU) 2021/664}.
\newblock Supplied project source, May 2024 edition.

\end{thebibliography}

\end{document}
